\chapter{Att kontrollera ett påstående}
\label{app:verifiera}

Det här materialet gör, som allt undervisningsmaterial, en rad
påståenden.  De flesta kan du pröva själv genom att köra programmen ---
det är hela poängen med att koden på bilderna är samma kod som körs.  Men
några påståenden går inte att pröva på det sättet:
\begin{enumerate}
  \item att Pythons listor beter sig som materialet beskriver, också i de
    delar som inte visas här (\cref{sec:listdok});
  \item att missuppfattningar om listindex förekommer hos nybörjare
    (\cref{sec:listor});
  \item att det är bra att inte upprepa sig i ett program --- principen
    \emph{DRY} --- och att den härrör från Hunt och Thomas
    (\cref{sec:sokning});
  \item att en funktion ska ha ett ansvar och att lösningen ska hållas
    enkel --- principerna \foreignlanguage{english}{SRP} och
    \foreignlanguage{english}{KISS} (\cref{sec:delsokning});
  \item de didaktiska påståenden som anteckningarna i marginalen gör om
    hur människor lär sig.
\end{enumerate}
Hur vet man om sådant stämmer?  Samma fråga möter dig varje gång du läser
en lärobok, en webbsida eller ett svar från en språkmodell --- och varje
gång du själv ska skriva en rapport.  Den här bilagan beskriver en
arbetsgång för att kontrollera ett påstående.
\Cref{tab:listor-pastaenden} visar var vart och ett av materialets
påståenden är belagt.

\begin{table}[htbp]
  \centering
  \caption{Modulens påståenden, som frågor, och var svaret finns.}
  \label{tab:listor-pastaenden}
  \begin{tabular}{@{}p{0.46\linewidth}p{0.44\linewidth}@{}}
    \toprule
    Fråga & Var det är belagt \\
    \midrule
    Beter sig Pythons listor som materialet beskriver?
      & Primärkälla: Pythons egen dokumentation, läst i fulltext
        2026-09-03. \\
    Förekommer missuppfattningar om listindex hos nybörjare?
      & Ja, enligt sammanfattningen av Veerasamy med flera; fulltexten
        har inte kunnat läsas härifrån, så påståendet är begränsat till
        att de förekommer. \\
    Är det bra att inte upprepa sig i ett program, och kommer principen
    \emph{DRY} från Hunt och Thomas?
      & Redan belagt i föreläsningen \emph{Algoritmiskt tänkande},
        bilaga~E: ja, med
        förbehåll. \\
    Härrör \foreignlanguage{english}{SRP} från Martin, och är den
    etablerad?
      & Redan belagt i föreläsningen \emph{Funktioner}, bilaga~B: ja,
        men vad som räknas som \emph{ett} ansvar är en
        bedömningsfråga. \\
    Är enkelhet ett etablerat designvärde, och kommer
    \foreignlanguage{english}{KISS} från Kelly Johnson?
      & Redan belagt i föreläsningen \emph{Funktioner}, bilaga~C: ja på
        den första frågan, nej på den andra. \\
    Måste kontrast upplevas samtidigt för att en aspekt ska kunna
    urskiljas?
      & Primärkälla: Marton, \emph{Necessary Conditions of Learning},
        s.~45--47. \\
    Varifrån kommer materialets index- och identitetskontraster?
      & Primärkälla: kursens eget manuskript om missuppfattningar, läst
        i fulltext 2026-09-16. \\
    \bottomrule
  \end{tabular}
\end{table}

\section{Gör påståendet till en fråga}

Ett påstående går att kontrollera först när det är tydligt vad som skulle
kunna visa att det är \emph{fel}.  Formulera det därför som en fråga med
ett svar: \enquote{tar det längre tid att söka i en lista ju längre listan
är?} kan besvaras med ja eller nej, medan \enquote{listor är viktiga} inte
kan det --- ingen tänkbar källa skulle visa att det är fel, eftersom
\enquote{viktigt} är ett omdöme och inte ett faktum.  Ofta döljer sig
flera påståenden i ett: \enquote{stegvis förfining är en etablerad metod
som härrör från Wirth} --- ett påstående du mötte i föreläsningen
\emph{Algoritmiskt tänkande} --- är två frågor, om \emph{ursprung} och om
\emph{etablering}, och de kräver olika slags källor.

\section{Välj rätt sorts källa}

Olika frågor besvaras av olika källor.
\begin{description}
  \item[Uppslagsverk] för vad ett ord \emph{betyder}.
  \item[Primärkällan] för var en idé \emph{kommer ifrån}, och för hur ett
    system \emph{beter sig}.  Frågan om vad en av Pythons behållare gör
    besvaras av Pythons egen dokumentation, inte av en lärobok som
    återger den.
  \item[Översikter] (\foreignlanguage{english}{reviews}) för om något är
    \emph{etablerat}, eller för vad forskningen \emph{sammantaget} säger.
  \item[Enskilda studier] för ett specifikt resultat --- med förbehållet
    att en enskild studie bara visar att något observerats en gång.
\end{description}
Var man hittar dem: universitetsbibliotekets söktjänst, och databaser som
Scopus, Web of Science, IEEE Xplore, ACM Digital Library, DBLP (datalogi)
och OpenAlex (öppen, och trots namnet inte detsamma som öppen tillgång).
Google Scholar täcker brett men rankar ogenomskinligt.

\section{Sök så att du kan ha fel}

Den som söker på namn hon redan känner till hittar det hon väntar sig.
Sök därför \emph{ämnesdrivet} --- på begrepp, inte författare.  Några
praktiska regler: håll frågorna korta, ställ samma frågor i flera
databaser, sök begreppets alla namn, och sök också efter
\emph{invändningar} genom att lägga till ord som \enquote{critique},
\enquote{limitations} eller \enquote{failure}.  Ett påstående som
överlevt en motsökning är värt mer än ett som aldrig prövats, och en
invändning som hittas blir ett \emph{förbehåll} i svaret, inte ett
nederlag.

\section{Läs i fulltext, och skriv ner hur}

En träff i en databas är inte ett belägg, och sammanfattningen
(\foreignlanguage{english}{abstract}) räcker inte heller: den säger vad
författarna vill ha sagt, inte vad de visat.  Läs hela texten och leta
upp \emph{det ställe} som styrker påståendet --- ett ordagrant citat med
sidnummer.  Räkna aldrig en källa du inte läst.

Skriv sedan ner hur du gick till väga, så att någon annan kan följa
spåret.  I det här materialet står den anteckningen som en kommentar
ovanför varje post i litteraturlistans källfil, med fasta fält:
\begin{description}
  \item[\texttt{CLAIM}] vilket påstående källan ska styrka;
  \item[\texttt{FOUND-VIA}] hur den hittades;
  \item[\texttt{PICKED}] varför just den, framför andra träffar;
  \item[\texttt{QUOTE}] det ordagranna citatet, med sida eller avsnitt;
  \item[\texttt{VERIFIED}] att fulltexten lästs, och var;
  \item[\texttt{COUNTER}] vad motsökningen gav;
  \item[\texttt{DATE}] när.
\end{description}
Det tar ett par minuter per källa, och det är det enda som gör det
möjligt att om ett år svara på frågan \enquote{hur vet du det?}

\section{Dra slutsatsen --- med förbehåll}
\label{sec:sok-slutsats}

Svaret skrivs ut tillsammans med det som talar emot.  Ett påstående om
hur snabb en operation är gäller till exempel en viss implementation och
ett genomsnittsfall, inte alla tänkbara fall --- och det förbehållet hör
till svaret.  Förbehållen är inte svagheter i svaret.  De \emph{är}
svaret, och de avgör hur påståendet får användas i materialet.

\section{En anmärkning om språkmodeller}

En språkmodell kan användas för att \emph{hitta} och \emph{sortera}
källor.  Lägg märke till arbetsfördelningen: varje källa som citeras ska
läsas och väljas av en människa.  Använd gärna språkmodeller för att
hitta och sortera; låt dem aldrig vara det sista ledet.  De kan göra
fel, men det är du själv som måste stå till svars för innehållet.

\section{Materialets påståenden och var de är belagda}
\label{sec:listor-pastaenden}

De flesta av den här modulens påståenden krävde ingen egen
litteratursökning: de är antingen frågor om vad Python gör --- och
besvaras då av Pythons egen dokumentation\autocite{PythonDataStructures},
som är primärkällan --- eller påståenden som redan är belagda någon
annanstans.  Ett påstående är hämtat ur en enskild studie, och det
redovisas nedan med sitt förbehåll.
Det vetenskapliga underlaget för att DRY är en etablerad princip som
härrör från Hunt och Thomas finns i föreläsningen \emph{Algoritmiskt
tänkande}, bilaga~E, med svaret \enquote{ja, med
förbehåll}\autocite{HuntThomas1999}.  Det vetenskapliga underlaget för
att \foreignlanguage{english}{SRP} härrör från Martin och är etablerad
--- med förbehållet att vad som räknas som \emph{ett} ansvar är en
bedömningsfråga --- finns i föreläsningen \emph{Funktioner}, bilaga~B,
och underlaget för att enkelhet är ett etablerat designvärde, medan
tillskrivningen av \foreignlanguage{english}{KISS} till Kelly Johnson
inte går att bekräfta, i bilaga~C i samma föreläsning.
Att kontrast måste upplevas samtidigt för att en aspekt ska
kunna urskiljas är hämtat från Martons \emph{Necessary Conditions of
Learning}\autocite[45--47]{Marton2015}.  Att återanvända ett belägg i
stället för att göra om sökningen är samma DRY-princip som materialet
lär ut, tillämpad på litteraturen.  Varifrån materialets egna
variationsmönster kommer --- indexkontrasten och identitetskontrasten ---
är inte ett påstående om världen utan en hänvisning, och den går till
kursens eget manuskript om
missuppfattningar\autocite{SooriBosk2026}, läst i fulltext.

Ett påstående är belagt med förbehåll.  Att missuppfattningar om
listreferenser och listindex förekommer hos nybörjare kommer från
Veerasamy med
fleras\autocite{KumarVeerasamy2016} studie av en tentamen i
introduktionsprogrammering.  Den artikeln har inte kunnat läsas i
fulltext härifrån --- den öppna kopian svarar med ett fel och
förlagsversionen ligger bakom betalvägg --- så materialet lutar sig bara
mot det som står i sammanfattningen, att missuppfattningar iakttogs
kring \enquote{list referencing} och inbyggda funktioner.  Hur
\emph{vanliga} de är påstår materialet därför ingenting om.  Det är ett
exempel på förbehållet i \cref{sec:sok-slutsats}: säg vad källan bär,
inte vad den skulle kunna bära.  Svaret på vart och ett av modulens
påståenden står sammanställt i \cref{tab:listor-pastaenden}.

